\documentclass[a4paper]{article}
\usepackage{amsfonts}
\usepackage{amsmath}
\usepackage{amssymb}
\usepackage{graphicx}     

\begin{document}

\noindent \large Auteur: Anna Voronovska \\
\noindent \large Studierichting: Informatica \\
\noindent \large Oefeningenreeks 1C nr. 31.2 \\

\bigskip 

\normalsize
Gegeven:\\
 
$ (-1 + i)^{10} $\\

\textbf{$z = -1 + i$ in poolvorm:}\\

De modulus:\\

$ r = \sqrt{(-1)^2 + 1^2} = \sqrt{1 + 1} = \sqrt{2} $\\

Het argument:\\

$ \tan \alpha = \dfrac{1}{-1} = -1 $;
$ \alpha = \dfrac{3\pi}{4} $\\

Dus de poolvorm is:\\

$\left( \sqrt{2} \operatorname{cis} \left(\dfrac{3\pi}{4}\right) \right)$\\


\textbf{Stelling van de Moivre:}\\

$\left(\sqrt{2} \operatorname{cis} \left(\dfrac{3\pi}{4}\right) \right)^{10} =
\left(\sqrt{2}\right)^{10} \operatorname{cis} \left(10 \cdot \dfrac{3\pi}{4} \right) =
32 \operatorname{cis} \left(\dfrac{30\pi}{4}\right) = 32 \operatorname{cis} \left(\dfrac{15\pi}{2}\right)$

Omdat argumenten herhalen na $ 2\pi $:\\

$ \dfrac{15\pi}{2} \mod 2\pi  = \dfrac{\pi}{2} $\\

Dus:\\

$ (-1 + i)^{10} = 32 \operatorname{cis} \left(\dfrac{\pi}{2}\right) $\\

Omdat $ \operatorname{cis} \left(\dfrac{\pi}{2}\right) = \operatorname{cis} 90^\circ = i $, krijgen wij:
$ (-1 + i)^{10} = 32i $\\

\textbf{Antwoord:}\\

$ (-1 + i)^{10} = 32i $\\

\begin{thebibliography}{999}
%\bibitem{a} Referentie
\end{thebibliography}

\end{document}
